\subsection{Integración y consumo de APIs: arquitectura REST y servicios externos}
\label{subsec:integracion_apis}

En este apartado se detalla el flujo de comunicación y consumo de interfaces de programación de aplicaciones (APIs) implementado en Colivi, abarcando tanto la interacción entre el cliente web y el backend propio mediante arquitectura REST, como la integración con servicios públicos externos de geolocalización. El diseño persigue dos metas esenciales de ingeniería: garantizar contratos de datos estrictamente tipados y desacoplados, y optimizar el rendimiento minimizando las llamadas a servicios externos sujetos a restricciones de cuota.

\subsubsection{Consumo e integración de la API REST interna (Frontend -- Backend)}

La comunicación entre el cliente desarrollado en React y los microservicios transaccionales de Spring Boot se estructura sobre un protocolo RESTful sin estado (\emph{stateless}). Para evitar discrepancias de datos y fallos en tiempo de ejecución, se aplicó un principio de modelado simétrico entre los \emph{Data Transfer Objects} (DTOs) de Java y las interfaces tipadas de TypeScript en el cliente:

\begin{itemize}
    \item \textbf{Gestión de peticiones y mutaciones reactivas}: Las consultas de lectura (\texttt{GET}) y las operaciones de mutación (\texttt{POST}, \texttt{PUT}, \texttt{DELETE}) se gestionan a través de ganchos personalizados (\emph{custom hooks}) que encapsulan las funciones de \textbf{TanStack Query} (\texttt{useQuery} y \texttt{useMutation}). Esta capa intermedia proporciona un manejo automático de los estados del ciclo de vida asíncrono (\texttt{isLoading}, \texttt{isError}, \texttt{data}), evitando la duplicación de código en los componentes visuales.
    \item \textbf{Propagación de credenciales y autorización}: Mediante el interceptor de Axios, cada solicitud hacia rutas protegidas viaja acompañada de la cabecera estándar \texttt{Authorization: Bearer <JWT>}. El backend valida la firma criptográfica en el filtro de seguridad antes de despachar la petición al controlador de dominio pertinente.
    \item \textbf{Políticas de invalidación y actualización optimista}: Para garantizar una experiencia fluida al usuario, las mutaciones desencadenan la invalidación reactiva de claves de consulta específicas (\texttt{queryClient.invalidateQueries}). Por ejemplo, al registrar la liquidación de una deuda compartida en el hogar, el libro mayor contable y los saldos cruzados de todos los compañeros se sincronizan automáticamente en el cliente sin requerir una recarga completa de la página.
\end{itemize}

\subsubsection{Consumo de APIs externas y geocodificación (OpenStreetMap / Nominatim)}

Una de las funcionalidades nucleares de Colivi radica en la exploración de inmuebles sobre un mapa cartográfico en tiempo real. No obstante, los anfitriones y propietarios introducen la ubicación de sus viviendas mediante direcciones postales en lenguaje natural (calle, número, código postal y ciudad), las cuales no pueden ser interpretadas directamente por el motor cartográfico de Leaflet sin antes convertirse en coordenadas espaciales $(\text{latitud}, \text{longitud})$.

Para resolver esta necesidad se integró la API pública de geocodificación \textbf{Nominatim}, provista por la fundación \textbf{OpenStreetMap}. Sin embargo, este servicio externo impone una estricta Política de Uso Aceptable (\emph{Acceptable Use Policy}):
\begin{itemize}
    \item Un límite de tasa estricto (*Rate Limiting*) de \textbf{máximo 1 petición por segundo}.
    \item Prohibición expresa de ejecutar consultas masivas o búsquedas repetitivas para una misma dirección.
    \item Requerimiento obligatorio de cabecera \texttt{User-Agent} identificativa para monitorizar el consumo del servicio.
\end{itemize}

Si la aplicación web consultara la API externa de Nominatim cada vez que un usuario explora el mapa o desplaza el zoom, el sistema colapsaría de inmediato al superar los límites de tasa, acarreando el bloqueo de la dirección IP del servidor y generando una latencia inasumible en la interfaz.

Para sortear esta limitación y garantizar la máxima eficiencia, se implementó una estrategia de \textbf{geocodificación asíncrona con persistencia local}:

\begin{enumerate}
    \item \textbf{Resolución única durante el alta del inmueble}: Cuando el anfitrión publica o actualiza un alojamiento, el backend captura la dirección estructurada y ejecuta una única llamada controlada al servicio de Nominatim.
    \item \textbf{Persistencia de coordenadas en base de datos}: La respuesta JSON de la API externa es procesada y las coordenadas resultantes se almacenan permanentemente en las columnas \texttt{latitude} y \texttt{longitude} de la tabla relacional \texttt{accommodations} de PostgreSQL.
    \item \textbf{Consultas locales indexadas}: En las búsquedas habituales del cliente web, la plataforma no realiza ninguna llamada a Nominatim; las viviendas se consultan directamente desde la base de datos local mediante índices numéricos sobre las coordenadas geográficas, sirviendo los datos con tiempos de respuesta inferiores a 20~ms y respetando rigurosamente los términos de uso del proveedor libre.
\end{enumerate}

\subsubsection{Verificación inicial de contratos y datos de prueba}

Con el fin de verificar la robustez de los contratos de la API REST y certificar visualmente la renderización de las coordenadas en el mapa cartográfico sin depender de inserciones manuales tediosas durante el desarrollo, se diseñó una inicialización automatizada de datos de prueba (\emph{seed data}). Esta se articula mediante una migración determinista en Flyway (\texttt{V999\_\_insert\_dev\_seed\_data.sql}) que preconfigura usuarios maestros y hogares de convivencia con cláusula \texttt{ON CONFLICT DO NOTHING}, complementada por un script en Python (\texttt{seed\_db.py}) que simula la interacción de anfitriones publicando inmuebles geolocalizados en los distritos universitarios más representativos de España (Madrid, Barcelona, Granada, Salamanca, Sevilla y Valencia).

